% Part: methods
% Chapter: induction
% Section: induction-on-N

\documentclass[../../../include/open-logic-section]{subfiles}

\begin{document}

\olfileid[id]{mth}{ind}{inN}

\olsection{Induksi pada~$\Nat$}

Dalam bentuknya yang paling sederhana, induksi merupakan teknik untuk
membuktikan hasil bagi semua bilangan asli. Teknik ini menggunakan fakta
bahwa, dengan mulai dari $0$ dan berulang kali menambahkan~$1$, kita pada
akhirnya mencapai setiap bilangan asli. Jadi, untuk membuktikan bahwa
sesuatu benar bagi setiap bilangan, kita dapat (1)~menetapkan bahwa hal itu
benar bagi $0$ dan (2)~menunjukkan bahwa setiap kali hal itu benar bagi suatu
bilangan~$n$, hal itu juga benar bagi bilangan berikutnya~$n+1$. Jika kita
menyingkat ``bilangan~$n$ memiliki sifat $P$'' menjadi $P(n)$ (dan
``bilangan~$k$ memiliki sifat $P$'' menjadi $P(k)$, dan seterusnya), maka
bukti induksi bahwa $P(n)$ bagi semua $n \in \Nat$ terdiri atas:
\begin{enumerate}
\item bukti bagi $P(0)$, dan
\item bukti bahwa, untuk sebarang~$k$, jika $P(k)$ maka $P(k+1)$.
\end{enumerate}
Untuk memperjelasnya sepenuhnya, andaikan kita memiliki (1) dan~(2). Maka
(1) memberi tahu kita bahwa $P(0)$ benar. Jika kita juga memiliki~(2), kita
khususnya mengetahui bahwa jika $P(0)$ maka $P(0+1)$, yakni $P(1)$. Hal ini
mengikuti pernyataan umum ``untuk sebarang~$k$, jika $P(k)$ maka $P(k+1)$''
dengan mengganti~$k$ dengan~$0$. Jadi, melalui modus ponens, kita memperoleh
$P(1)$. Dari (2) lagi, sekarang dengan mengganti~$k$ dengan~$1$, kita
memperoleh: jika $P(1)$ maka~$P(2)$. Karena kita baru saja menetapkan
$P(1)$, melalui modus ponens kita memperoleh~$P(2)$. Demikian seterusnya.
Untuk sebarang bilangan~$n$, setelah melakukan langkah ini sebanyak
$n$~kali, kita akhirnya tiba pada~$P(n)$. Jadi, (1) dan~(2) bersama-sama
menetapkan~$P(n)$ bagi sebarang $n \in \Nat$.

Mari kita lihat sebuah contoh. Andaikan kita ingin mengetahui berapa banyak
jumlah berbeda yang dapat dihasilkan oleh lemparan $n$ dadu. Walaupun
mungkin tampak konyol, mari kita mulai dengan $0$ dadu. Jika tidak memiliki
dadu, hanya ada satu jumlah yang mungkin ``dilempar'': tidak ada mata dadu
sama sekali, sehingga jumlahnya~$0$. Jadi, banyaknya hasil berbeda yang
mungkin adalah~$1$. Jika Anda hanya memiliki satu dadu, yakni $n=1$, ada
enam nilai yang mungkin, dari $1$ sampai~$6$. Dengan dua dadu, kita dapat
menghasilkan setiap jumlah dari $2$ sampai $12$, yakni $11$ kemungkinan.
Dengan tiga dadu, kita dapat menghasilkan setiap bilangan dari $3$ sampai
$18$, yakni $16$ kemungkinan berbeda. $1$, $6$, $11$, $16$: tampak seperti
suatu pola; mungkinkah jawabannya $5n+1$? Tentu saja, $5n+1$ adalah jumlah
maksimum kemungkinan, sebab hanya ada $5n+1$ bilangan di antara $n$, nilai
terendah yang dapat dihasilkan dengan $n$ dadu (semuanya menunjukkan $1$),
dan $6n$, nilai tertinggi yang dapat dihasilkan (semuanya menunjukkan $6$).

\begin{thm}
  Dengan $n$ dadu, semua $5n+1$ nilai yang mungkin di antara $n$ dan $6n$
  dapat dihasilkan.
\end{thm}

\begin{proof}
Misalkan $P(n)$ adalah klaim: ``Dengan $n$~dadu, setiap bilangan di antara
$n$ dan~$6n$ dapat dihasilkan.'' Untuk menggunakan induksi, kita buktikan:
\begin{enumerate}
\item \emph{dasar induksi} $P(0)$, yakni dengan tanpa dadu, kita dapat
  menghasilkan satu-satunya bilangan di antara $0$ dan $0$;
\item \emph{langkah induksi}: untuk semua $k$, jika $P(k)$ maka~$P(k+1)$.
\end{enumerate}

(1) benar karena tanpa dadu tidak ada mata dadu yang dijumlahkan, sehingga
satu-satunya jumlah yang dihasilkan adalah~$0$.

Untuk membuktikan (2), kita mengasumsikan anteseden kondisional tersebut,
yakni $P(k)$. Asumsi ini disebut \emph{hipotesis induksi}. Kita
menggunakannya untuk membuktikan~$P(k+1)$. Bagian yang sulit ialah menemukan
cara memikirkan nilai-nilai yang mungkin dari lemparan $k+1$ dadu berdasarkan
nilai-nilai yang mungkin dari lemparan $k$~dadu ditambah lemparan dadu
tambahan yang ke-$(k+1)$---tetapi inilah yang harus dilakukan jika kita
hendak menggunakan hipotesis induksi.

Hipotesis induksi menyatakan bahwa kita dapat memperoleh setiap bilangan di
antara $k$ dan~$6k$ dengan menggunakan $k$~dadu. Jika dadu ke-$(k+1)$
menunjukkan~$1$, hal ini menambahkan $1$ pada jumlah keseluruhan. Jadi, kita
dapat menghasilkan setiap nilai di antara $k+1$ dan $6k+1$ dengan melempar
$k$~dadu lalu memperoleh~$1$ pada dadu ke-$(k+1)$. Apa yang tersisa?
Nilai-nilai $6k+2$ sampai $6k+6$. Kita dapat memperolehnya dengan membuat
semua $k$ dadu pertama menunjukkan~$6$, lalu memperoleh bilangan antara $2$
dan $6$ pada dadu ke-$(k+1)$. Secara keseluruhan, ini berarti bahwa dengan
$k+1$ dadu kita dapat menghasilkan setiap bilangan di antara $k+1$ dan
$6(k+1)$; jadi, kita telah membuktikan~$P(k+1)$ dengan menggunakan asumsi
$P(k)$, yaitu hipotesis induksi.
\end{proof}

Kita sangat sering menggunakan induksi ketika hendak membuktikan sesuatu
mengenai suatu deret objek (bilangan, himpunan, dan sebagainya) yang dengan
sendirinya didefinisikan ``secara induktif'', yakni dengan mendefinisikan
objek ke-$(n+1)$ berdasarkan objek ke-$n$. Misalnya, kita dapat
mendefinisikan jumlah~$s_n$ dari bilangan-bilangan asli sampai~$n$ dengan
\begin{align*}
  s_0 & = 0\\
  s_{n+1} & = s_n + (n+1)
\end{align*}
Definisi ini memberikan:
\begin{align*}
  s_0 & = 0,\\
  s_1 & = s_0 + 1 && = 1,\\
  s_2 & = s_1 + 2 && = 1 + 2 = 3\\
  s_3 & = s_2 + 3 && = 1 + 2 + 3 = 6, \text{ dan seterusnya.}
\end{align*}
Sekarang kita dapat membuktikan dengan induksi bahwa $s_n = n(n+1)/2$.

\begin{prop}
  $s_n = n(n+1)/2$.
\end{prop}

\begin{proof}
  Kita harus membuktikan (1) bahwa $s_0 = 0\cdot(0 + 1)/2$ dan (2) bahwa
  jika $s_k = k(k+1)/2$, maka $s_{k+1} = (k+1)(k+2)/2$. (1) sudah jelas.
  Untuk membuktikan (2), kita mengasumsikan hipotesis induksi: $s_k =
  k(k+1)/2$. Dengan menggunakannya, kita harus menunjukkan bahwa $s_{k+1}
  = (k+1)(k+2)/2$.

  Berapakah $s_{k+1}$? Menurut definisi, $s_{k+1} = s_k + (k+1)$. Menurut
  hipotesis induksi, $s_k = k(k+1)/2$. Kita dapat menyubstitusikannya ke
  persamaan sebelumnya, lalu hanya memerlukan sedikit aritmetika pecahan:
  \begin{align*}
    s_{k+1} & = \frac{k(k+1)}{2} + (k+1) \\
    & = \frac{k(k+1)}{2} + \frac{2(k+1)}{2} \\
    & = \frac{k(k+1) + 2(k+1)}{2} \\
    & = \frac{(k+2)(k+1)}{2}.
  \end{align*}
\end{proof}

Pelajaran pentingnya ialah bahwa jika Anda membuktikan sesuatu mengenai
suatu barisan $a_n$ yang didefinisikan secara induktif, induksi merupakan
cara yang sudah semestinya dipilih. Bahkan jika tidak demikian (seperti
dalam kasus berbagai kemungkinan hasil lemparan dadu), Anda dapat
menggunakan induksi apabila dapat menghubungkan kasus untuk~$k+1$ dengan
kasus untuk~$k$.

\end{document}
